
\documentclass[a4paper,11pt]{article}
\usepackage{fontspec}
\setmainfont{Noto Sans}
\usepackage[french]{babel}
\usepackage{microtype}
\usepackage[top=9mm,bottom=8mm,left=11mm,right=11mm]{geometry}
\usepackage{xcolor}
\usepackage{tikz}
\usetikzlibrary{arrows.meta,calc,positioning}
\usepackage[most]{tcolorbox}
\usepackage{ulem}
\pagestyle{empty}

\definecolor{Green}{HTML}{319B58}
\definecolor{GreenDark}{HTML}{17613A}
\definecolor{GreenLight}{HTML}{EAF6EC}
\definecolor{GreenLine}{HTML}{8FC9A3}
\definecolor{GreenBorder}{HTML}{B9D3C3}
\definecolor{Orange}{HTML}{D8842F}
\definecolor{OrangeDark}{HTML}{9A5716}
\definecolor{OrangeLight}{HTML}{FFF1DF}
\definecolor{OrangeLine}{HTML}{E7B77E}
\definecolor{OrangeBorder}{HTML}{E7C7A6}
\definecolor{Violet}{HTML}{7A5AA6}
\definecolor{VioletDark}{HTML}{563A7A}
\definecolor{VioletLight}{HTML}{F2ECF8}
\definecolor{VioletLine}{HTML}{B9A5D0}
\definecolor{VioletBorder}{HTML}{D3C4E2}
\definecolor{Ink}{HTML}{172A32}
\definecolor{Grey}{HTML}{64757D}
\definecolor{Red}{HTML}{D52F35}
\definecolor{RedLight}{HTML}{FCEBED}

\tcbset{
 enhanced,boxrule=.6pt,arc=2.8mm,
 left=3.7mm,right=3.7mm,top=1.75mm,bottom=1.75mm,
 boxsep=1mm,coltext=Ink
}
\newtcolorbox{greenbox}[1][]{colback=GreenLight,colframe=GreenBorder,#1}
\newtcolorbox{orangebox}[1][]{colback=OrangeLight,colframe=OrangeBorder,#1}
\newtcolorbox{violetbox}[1][]{colback=VioletLight,colframe=VioletBorder,#1}
\newtcolorbox{warnbox}[1][]{colback=RedLight,colframe=Red,
 left=2.7mm,right=2.7mm,top=1.25mm,bottom=1.25mm,#1}

\begin{document}
\begin{tikzpicture}[remember picture,overlay]
\draw[draw=OrangeBorder,line width=.9pt,rounded corners=1.7mm]
 ([xshift=1mm,yshift=-1mm]current page.north west)
 -- ([xshift=-1mm,yshift=-1mm]current page.north east)
 -- ([xshift=-1mm,yshift=1mm]current page.south east)
 -- ([xshift=1mm,yshift=1mm]current page.south west)--cycle;
\fill[Orange] ([xshift=1mm,yshift=-1mm]current page.north west)
 rectangle ([xshift=5.2mm,yshift=1mm]current page.south west);
\end{tikzpicture}

\vspace*{1mm}\hspace*{6mm}
{\color{OrangeDark}\fontsize{25}{27}\selectfont\bfseries FUTUR PROCHE}

\vspace{0.5mm}\hspace*{6mm}
{\color{Grey}\large Le futur proche}

\vspace{2.5mm}\hspace*{6mm}
\begin{minipage}{0.91\textwidth}

\begin{orangebox}
\textbf{\color{OrangeDark}CONSTRUCTION}\\[1mm]
\centering
{\Large \textbf{ALLER AU PRÉSENT + INFINITIF}}\\[1mm]
\small
je vais \quad tu vas \quad il/elle va \quad nous allons \quad vous allez \quad ils/elles vont
\end{orangebox}

\vspace{2.5mm}
\begin{center}
\begin{tikzpicture}[>=Latex]
\draw[Ink!65,line width=.8pt] (0,0)--(15.1,0);
\node[below=2mm] at (2.0,0) {\small PASSÉ};
\node[below=2mm] at (7.5,0) {\small PRÉSENT};
\node[below=2mm] at (13.5,0) {\small FUTUR};
\draw[Orange,line width=2.5pt,->] (8.0,0)--(13.0,0);
\draw[OrangeDark,line width=1.1pt,->] (7.5,1.0)--(7.5,.18);
\node[above=3.8mm] at (7.5,1.0) {\bfseries\color{OrangeDark}MAINTENANT};
\node[above=3mm,align=center,text width=6.2cm] at (10.5,0)
{\small\bfseries futur proche : action envisagée ou imminente};
\end{tikzpicture}
\end{center}

\vspace{1.5mm}
\begin{orangebox}
\textbf{\color{OrangeDark}À QUOI SERT-IL ?}\\[1mm]
\begin{itemize}\setlength{\itemsep}{0.4mm}
\item parler d'une action qui va se produire \textbf{bientôt} ;
\item exprimer une intention ou un projet déjà envisagé ;
\item annoncer une conséquence qui paraît \textbf{imminente} ;
\item faire une prédiction fondée sur des signes présents.
\end{itemize}
\end{orangebox}

\vspace{2mm}
\begin{orangebox}
\textbf{\color{OrangeDark}EXEMPLES}\\[1mm]
\textbf{Je vais partir dans quelques minutes.}\\
\textbf{Nous allons commencer le travail.}\\
\textbf{Regarde ces nuages : il va pleuvoir.}\\
\textbf{Elle va appeler le médecin ce soir.}
\end{orangebox}

\vspace{2mm}
\begin{minipage}{0.485\textwidth}
\begin{orangebox}[height=3.15cm]
\textbf{\color{OrangeDark}FUTUR PROCHE / FUTUR SIMPLE}\\[1mm]
Le futur proche met souvent l'accent sur la \textbf{proximité}, l'intention ou les signes présents.\\[1mm]
\textit{Je vais partir.}\\
\textit{Je partirai demain.}
\end{orangebox}
\end{minipage}\hfill
\begin{minipage}{0.485\textwidth}
\begin{warnbox}[height=3.15cm]
\textbf{\color{Red}ATTENTION}\\[1mm]
\textbf{Aller} est ici un \textbf{auxiliaire} au présent. L'action principale reste à l'infinitif.\\[1mm]
\textit{Je vais \textbf{manger}.}\\
\textcolor{Grey}{et non : \sout{je vais mangé}}
\end{warnbox}
\end{minipage}

\vspace{2mm}
\begin{orangebox}
\textbf{\color{OrangeDark}À RETENIR}\\[1mm]
Le futur proche est très courant dans la langue parlée. Il permet de présenter une action future comme \textbf{proche, envisagée ou déjà engagée}.
\end{orangebox}

\vspace{1.5mm}\centering
{\small\color{Grey}\textbf{Repères fréquents :} bientôt, dans un instant, tout à l'heure, ce soir, dans quelques minutes, prochainement...}

\end{minipage}
\end{document}
